\DocumentMetadata{
  pdfversion=2.0,
  pdfstandard=ua-2,
  lang=en-US,
  tagging=on,
  tagging-setup={math/setup=mathml-SE,tabsorder=structure}
}
\documentclass[11pt,letterpaper]{article}
\usepackage[margin=0.68in,includefoot]{geometry}
\usepackage{newcomputermodern}
\usepackage{amsmath,amscd,mathtools}
\usepackage{tikz,adjustbox}
\providecommand{\square}{\mdlgwhtsquare}
\usepackage[hidelinks]{hyperref}
\hypersetup{
  pdftitle={Numerical Analysis First-Year Exam, January 2013},
  pdfauthor={University of Florida Department of Mathematics},
  pdfsubject={Graduate mathematics examination},
  pdfdisplaydoctitle=true,
  bookmarksopen=true
}
\setcounter{secnumdepth}{0}
\setlength{\parindent}{0pt}
\setlength{\parskip}{0.28em}
\setlength{\emergencystretch}{3em}
\setlength{\leftmargini}{2.15em}
\setlength{\leftmarginii}{2.65em}
\setlength{\leftmarginiii}{3em}
\pagestyle{empty}
\raggedbottom
\allowdisplaybreaks
\NewTaggingSocketPlug{block/list/label}{examproblem}{%
  \tagstructbegin{tag=Lbl}\tagstructend%
  \tagstructbegin{tag=\LBody}%
  \tagstructbegin{tag=\ProblemHeadingTag,title-o={\ProblemHeadingText}}%
  \tagmcbegin{tag=\ProblemHeadingTag,actualtext={\ProblemHeadingText}}%
  #2%
  \tagmcend%
  \tagstructend%
}
\newenvironment{examproblems}[2]{%
  \def\ProblemHeadingTag{#2}%
  \AssignTaggingSocketPlug{block/list/label}{examproblem}%
  \begin{enumerate}%
  \setcounter{enumi}{#1}%
  \renewcommand{\labelenumi}{\textbf{\theenumi.}}%
  \setlength{\topsep}{0.35em}%
  \setlength{\partopsep}{0pt}%
  \setlength{\itemsep}{0.48em}%
  \setlength{\parsep}{0.18em}%
}{\end{enumerate}}
\newenvironment{examsubparts}{%
  \AssignTaggingSocketPlug{block/list/label}{default}%
  \begin{enumerate}%
  \setlength{\topsep}{0.18em}%
  \setlength{\partopsep}{0pt}%
  \setlength{\itemsep}{0.16em}%
  \setlength{\parsep}{0.1em}%
}{\end{enumerate}}
\newenvironment{examinstructions}{%
  \par\begingroup\itshape
}{%
  \par\endgroup\vspace{0.18em}
}
\makeatletter
\renewcommand\subsection{\@startsection{subsection}{2}{0pt}%
  {-1.25ex plus -.25ex minus -.1ex}%
  {0.55ex plus .1ex}%
  {\normalfont\large\bfseries}}
\renewcommand{\@maketitle}{%
  \newpage\null\vskip -1em%
  {\footnotesize\bfseries
    \href{https://www.ufl.edu/}{University of Florida}, \href{https://math.ufl.edu/}{Department of Mathematics}\par}%
  \vskip -0.5em%
  \tagstructbegin{tag=H1,title={Numerical Analysis First-Year Exam, January 2013}}%
  \tagpdfparaOff%
  \tagmcbegin{tag=H1}%
    {\LARGE\bfseries\href{https://gma.math.ufl.edu/exams/numerical-analysis-first-year/}{Numerical Analysis First-Year Exam}, January 2013\par}%
  \tagmcend%
  \tagpdfparaOn%
  \tagstructend%
  \vskip 0.25em%
  \tagmcbegin{artifact}%
    \hrule height 1pt%
  \tagmcend%
  \vskip 1em%
}
\makeatother
\title{Numerical Analysis First-Year Exam, January 2013}
\author{}
\date{}
\begin{document}
\maketitle
\thispagestyle{empty}
\begin{examinstructions}
Answer any 4 questions.
\end{examinstructions}
\begin{examproblems}{0}{H2}
\def\ProblemHeadingText{Problem 1.}
\item
Suppose that \(A \in \mathbb{C}^{n\times n}\) is Hermitian with \(\lambda_1\) and \(\lambda_n\) the smallest and largest eigenvalues of~\(A\) respectively. Show that
\[
\lambda_1\lVert x\rVert^2 \leq x^*Ax \leq \lambda_n\lVert x\rVert^2
\]
 for all \(x \in \mathbb{C}^n\).
\def\ProblemHeadingText{Problem 2.}
\item
This problem has three parts.
\begin{examsubparts}
\item[\textnormal{(a)}]
Suppose~\(p\) and \(q \in \mathbb{R}\) with~\(p\) and~\(q\) positive and \(p^{-1}+q^{-1}=1\). Show that for any matrix \(A \in \mathbb{C}^{m\times n}\), we have \(\lVert A\rVert_p=\lVert A^*\rVert_q\), where \(A^*\) is the conjugate transpose of~\(A\).
\item[\textnormal{(b)}]
Prove that
\[
\lVert A\rVert_2^2 \leq \lVert A\rVert_p\lVert A\rVert_q
\]
 for any \(A \in \mathbb{C}^{m\times n}\) and any positive~\(p\) and \(q \in \mathbb{R}\) with \(p^{-1}+q^{-1}=1\).
\item[\textnormal{(c)}]
Prove that for any \(p\geq 1\) and any diagonal matrix \(D \in \mathbb{C}^{n\times n}\), we have
\[
\lVert D\rVert_p=\max\{\lvert d_{ii}\rvert:1\leq i\leq n\}.
\]
\end{examsubparts}
\def\ProblemHeadingText{Problem 3.}
\item
This problem has two parts.
\begin{examsubparts}
\item[\textnormal{(a)}]
Show that the eigenvalues of a projector are either~\(0\) or~\(1\).
\item[\textnormal{(b)}]
Show that a projector~\(P\) is orthogonal if and only if \(P=P^*\).
\end{examsubparts}
\def\ProblemHeadingText{Problem 4.}
\item
Show that the element of largest magnitude in a Hermitian, positive definite matrix lies on the diagonal.
\def\ProblemHeadingText{Problem 5.}
\item
Suppose that \(A \in \mathbb{C}^{n\times n}\) is an invertible matrix, \(u \in \mathbb{C}^n\), and \(v \in \mathbb{C}^n\).
\begin{examsubparts}
\item[\textnormal{(a)}]
Show that
\[
\det(A-uv^*)=(1-v^*A^{-1}u)\det A.
\]
\item[\textnormal{(b)}]
Suppose that~\(A\) is a real diagonal matrix and \(u=v\) with \(u_i\neq 0\) for each~\(i\). Show that the eigenvalues of \(A-uu^*\) are the roots of
\[
\sum_{i=1}^n \frac{\lvert u_i\rvert^2}{a_{ii}-\lambda}=1.
\]
\end{examsubparts}
\end{examproblems}
\end{document}
